% Bozza per la call con il Prof. Rivolta — 24 settembre
% Proposte per i Capitoli 4-5, non ancora in main.tex

\documentclass[a4paper,12pt]{article}
\usepackage[utf8]{inputenc}
\usepackage[T1]{fontenc}
\usepackage[italian]{babel}
\usepackage[a4paper,margin=2.5cm]{geometry}
\usepackage{amssymb,amsmath}
\usepackage{booktabs}
\usepackage{csquotes}
\usepackage[hidelinks]{hyperref}
\usepackage[backend=biber, citestyle=verbose-ibid, bibstyle=numeric, sorting=none, giveninits=true]{biblatex}
\addbibresource{bibliografia.bib}

\title{Proposte di revisione dei Capitoli 4 e 5}
\author{Daniele Intra}
\date{24 settembre 2026}

\begin{document}
\maketitle

% A — Sez. 4.3: loss + EoT, L1/L2, soglia, (0,1), replicabilità
\section{Funzione obiettivo e \textit{Expectation over Transformation}}
\label{sec:obiettivo}

\textbf{Parametrizzazione della perturbazione.}
La perturbazione non è ottimizzata direttamente nello spazio dei pixel, ma
attraverso una variabile libera $\theta \in \mathbb{R}^{3 \times 200 \times 200}$:
\begin{equation}
\label{eq:param_sigmoide}
P = \sigma(\theta), \qquad \theta_0 \sim \mathcal{N}(0,\, 0.1^2)
\end{equation}
dove $\sigma$ è la funzione logistica applicata elemento per elemento. La
riparametrizzazione garantisce $P \in (0,1)$ senza operazioni di troncamento, il
cui gradiente è nullo al di fuori dell'intervallo ammesso; è equivalente alla
trasformazione tramite tangente iperbolica proposta da Carlini e
Wagner\footcite{carlini_2017}, poiché $\sigma(2w) = \tfrac{1}{2}(\tanh w + 1)$.

\textbf{Composizione.}
Sia $x$ un fotogramma dell'insieme di addestramento e $B(x)$ l'insieme dei
riquadri di persona validi (al più tre, i più estesi). Per un riquadro di
larghezza $w$ e altezza $h$, la perturbazione occupa una regione di
$0.5\,w \times 0.4\,h$ centrata orizzontalmente e posta a $0.33\,h$ dal bordo
superiore. Indicando con $M_t$ la maschera di opacità trasformata da $t$,
l'immagine composta è
\begin{equation}
\label{eq:composizione}
A(x, P, t) = x \odot (1 - M_t) + t(P) \odot M_t .
\end{equation}

\textbf{Funzione obiettivo.}
L'avversario vuole la mancata rilevazione del soggetto, evento non
differenziabile. Si minimizza quindi una funzione obiettivo surrogata,
monotona crescente nella confidenza residua del rilevatore sulla classe
persona. Siano $t_1, \dots, t_N$, con $N = 16$, realizzazioni indipendenti della
distribuzione di trasformazioni $\mathcal{T}$ (Tabella~\ref{tab:eot_parametri}),
$\mathcal{K}_i$ l'insieme delle $K = 20$ celle più confidenti, entro i riquadri
di persona, nella realizzazione $i$, e $\omega_j \in [0,1]$ il peso della cella
$j$ in funzione dell'altezza del riquadro di appartenenza
(Sezione~4.4). La confidenza aggregata è
\begin{equation}
\label{eq:conf_aggregata}
\bar{c}(x, P) =
\frac{\sum_{i=1}^{N} \sum_{j \in \mathcal{K}_i} \omega_j \,
      \sigma\big(f_j(A(x, P, t_i))\big)}
     {\sum_{i=1}^{N} \sum_{j \in \mathcal{K}_i} \omega_j}
\end{equation}
dove $f_j$ è il punteggio di classe persona restituito dal rilevatore sulla
cella $j$, prima della soppressione dei non-massimi. La funzione obiettivo,
valutata su quattro fotogrammi per aggiornamento, è
\begin{equation}
\label{eq:obiettivo}
\mathcal{J}(P) = \sum_{m=1}^{4} -\ln\!\big(1 - \bar{c}(x_m, P) + \varepsilon\big)
\;+\; \lambda_{TV}\, \mathcal{L}_{TV}(P),
\qquad \varepsilon = 10^{-6},\; \lambda_{TV} = 0.1 .
\end{equation}
Il termine logaritmico è strettamente crescente in $\bar{c}$ su tutto
l'intervallo $(0,1)$ e il suo gradiente non si annulla al decrescere della
confidenza.

\textbf{Relazione con l'\textit{Expectation over Transformation}.}
L'EoT di Athalye et al.\footcite{athalye_2018_eot} ottimizza il valore atteso
della perdita sulla distribuzione delle trasformazioni,
$\mathbb{E}_{t \sim \mathcal{T}}[\mathcal{L}(t)]$. In
\eqref{eq:conf_aggregata}--\eqref{eq:obiettivo} il valore atteso è invece
stimato sulla confidenza, prima della trasformazione logaritmica:
$-\ln(1 - \mathbb{E}_t[c])$. Poiché $-\ln(1-c)$ è convessa, per la
disuguaglianza di Jensen la quantità ottimizzata è un limite inferiore di
quella di Athalye: le realizzazioni in cui la perturbazione è meno efficace
pesano in misura uguale alle altre, anziché in misura maggiore.

\begin{table}[htbp]
\centering
\begin{tabular}{lll}
\toprule
\textbf{Componente} & \textbf{Distribuzione} & \textbf{Condizione simulata} \\
\midrule
Guadagno di contrasto & $\mathcal{U}(0.6,\,1.4)$ & risposta del sensore, esposizione \\
Offset di luminosità & $\mathcal{U}(-0.2,\,0.2)$ & illuminazione della scena \\
Rumore additivo, per pixel & $\mathcal{N}(0,\,0.05^2)$ & compressione, rumore di sensore \\
Rotazione & $\mathcal{U}(-15^{\circ},\,15^{\circ})$ & orientamento del soggetto \\
Fattore di scala & $\mathcal{U}(0.4,\,0.9)$ & distanza del velivolo \\
Rapporto d'aspetto & $\mathcal{U}(0.7,\,1.3)$ & inclinazione della ripresa \\
Traslazione, per asse & $\mathcal{U}(-0.1,\,0.1)$ & collocazione imprecisa sul corpo \\
\bottomrule
\end{tabular}
\caption{Componenti della distribuzione $\mathcal{T}$, estratte
indipendentemente per ciascuna realizzazione. Le componenti geometriche sono
composte in un'unica trasformazione affine con ricampionamento bilineare.}
\label{tab:eot_parametri}
\end{table}

\begin{table}[htbp]
\centering
\begin{tabular}{ll}
\toprule
\textbf{Parametro} & \textbf{Valore} \\
\midrule
Ottimizzatore & Adam su $\theta$, tasso di apprendimento iniziale $0.01$ \\
Programma del tasso & coseno fino a $0.001$ sull'intera durata \\
Aggiornamenti & $1250$ (VisDrone), $2000$ (Okutama); 4 fotogrammi ciascuno \\
Normalizzazione del gradiente & norma massima $1.0$ \\
Arresto anticipato & media mobile su $100$ valori, pazienza $1000$ \\
Selezione della perturbazione & minimo della media mobile della funzione obiettivo \\
\bottomrule
\end{tabular}
\caption{Parametri di ottimizzazione.}
\label{tab:param_ottimizzazione}
\end{table}

\textbf{Regolarizzazione.}
$\mathcal{L}_{TV}(P) = \operatorname{mean}|P_{:,:,w+1} - P_{:,:,w}| +
\operatorname{mean}|P_{:,h+1,:} - P_{:,h,:}|$, nella forma anisotropa in norma
$\ell_1$ adottata da Thys et al.\footcite{thys_2019}. Rispetto a una
penalizzazione quadratica, che attenua i salti ampi e tollera le piccole
variazioni producendo superfici sfumate, la norma $\ell_1$ spinge a zero le
differenze fra pixel adiacenti su regioni estese e ammette poche transizioni
nette: il risultato sono campiture uniformi, più fedelmente riproducibili su
supporto stampato.


% B — Sez. 4.1: automa x_{n+1}=f(x_n), distribuzioni, parametro estrazione, YOLO
\section{Il livello di simulazione come macchina a stati}
\label{sec:automa}

Il livello di simulazione non contiene il rilevatore. YOLOv8 è eseguito una
sola volta, in modalità non interattiva, sull'intero insieme di validazione,
prima dell'avvio della simulazione; richiamarlo a ogni passo per ogni entità
imporrebbe un numero di inferenze pari al prodotto fra entità e iterazioni.
Il simulatore riceve, attraverso il file di interfaccia, i quattro valori
$R_1^{\text{pre}}$, $R_1^{\text{post}}$, $R_2^{\text{pre}}$, $R_2^{\text{post}}$
e li impiega come parametri di un'estrazione di Bernoulli.

\textbf{Componenti.}

\begin{table}[htbp]
\centering
\begin{tabular}{lll}
\toprule
\textbf{Componente} & \textbf{Ingresso} & \textbf{Uscita} \\
\midrule
\texttt{Environment} & passo temporale & posizioni, entità visibili, civili prossimi \\
\texttt{VisionAgentStat} & entità, $R_1$, $R_2$ & esito di rilevamento, confidenza \\
\texttt{FusionAgent} & tre canali di evidenza & punteggio di minaccia, banda di variabilità \\
\texttt{DecisionAgent} & punteggio, banda, civili prossimi & azione \\
\bottomrule
\end{tabular}
\caption{Componenti del livello di simulazione e relative interfacce.}
\label{tab:componenti_sim}
\end{table}

\textbf{Formalizzazione.} Il sistema evolve a tempo discreto secondo
\begin{equation}
\label{eq:automa}
x_{n+1} = f(x_n, \xi_n), \qquad a_{e,n} = g(x_n, \xi'_n),
\end{equation}
dove $\xi_n, \xi'_n$ raccolgono le estrazioni casuali del passo $n$ e
$a_{e,n}$ è l'azione per l'entità $e$. Lo spazio degli stati non è finito,
poiché le componenti di $x_n$ assumono valori continui; è finito l'insieme
delle azioni (Tabella~\ref{tab:azioni}). Le azioni non fanno parte dello
stato: sono ricalcolate a ogni passo e non dipendono dall'azione precedente.

\textbf{Stato.} Per ogni entità $e$, $x_{e,n} = (p_{e,n},\, H_{e,n},\,
o_{e,n},\, h_{e,n})$: posizione intera in una griglia $100 \times 100$; ultime
$20$ posizioni; profilo OSINT $o = (b, g, s)$ con $b$ presenza della targa in
lista di sorveglianza, $g$ rischio geografico, $s$ numero di corrispondenze
sui canali social; ultimi $20$ punteggi di minaccia. Sono fissi per l'intera
simulazione il ruolo dell'entità, lo scenario e i quattro valori del file di
interfaccia.

\textbf{Inizializzazione.}

\begin{table}[htbp]
\centering
\begin{tabular}{lll}
\toprule
\textbf{Grandezza} & \textbf{Bersaglio} & \textbf{Civile} \\
\midrule
Numero di entità & $3$ & $30$ \\
Posizione & $\mathcal{U}\{0,\dots,99\}^2$ & $\mathcal{U}\{0,\dots,99\}^2$ \\
Targa in lista, $b$ & $\text{Bernoulli}(0.75)$ & $\text{Bernoulli}(0.01)$ \\
Rischio geografico, $g$ & $\text{Beta}(8,\,2)$ & $\text{Beta}(2,\,8)$ \\
Corrispondenze social, $s$ & $\text{Poisson}(12)$ & $\text{Poisson}(1)$ \\
\bottomrule
\end{tabular}
\caption{Distribuzioni di inizializzazione delle entità. Il velivolo parte dal
centro della griglia, quota $10$ m.}
\label{tab:init_sim}
\end{table}

\textbf{Transizione $f$.} A ogni passo, nell'ordine:
(i) ogni entità si sposta di $\delta_x, \delta_y \sim \mathcal{U}\{-k,\dots,k\}$
indipendenti, con $k = 2$ per i bersagli e $k = 1$ per i civili, entro i bordi
della griglia; (ii) il velivolo si sposta allo stesso modo con $k = 17$
($61$ km/h, assumendo un passo di $1$ s e celle di $1$ m); (iii) negli scenari
con attacco al canale OSINT, a ogni passo il profilo dei bersagli è portato a
$b = 0$, $s = 0$, $g \leftarrow \max(0,\, g - u)$ con
$u \sim \mathcal{U}(0.4,\, 0.6)$, e quello dei civili a $b = 1$, $g = 0.95$,
$s = 15$; (iv) la storia dei punteggi è aggiornata.

\textbf{Uscita $g$.} Per ciascuna entità, a ogni passo:

\emph{Canale visivo.} Esito $\delta \sim \text{Bernoulli}(q)$, con
$q = R_1^{\text{post}}$ per un bersaglio con contromisura attiva,
$q = R_1^{\text{pre}}$ per un bersaglio senza, $q = 1 - R_2$ per un civile.
Confidenza $c_v = \operatorname{clip}(u + \eta,\, 0,\, 1)$ con
$u \sim \mathcal{U}(0.50,\, 0.95)$ se $\delta = 1$,
$u \sim \mathcal{U}(0.05,\, 0.50)$ se $\delta = 0$, e
$\eta \sim \mathcal{N}(0,\, 0.02^2)$. Segnale $v = c_v$ se $\delta = 1$,
$v = 0.5\, c_v$ altrimenti. La distanza non interviene: $R_1$ è misurato su
fotogrammi reali a distanze eterogenee e ne incorpora già l'effetto.

\emph{Canale OSINT.} Classificatore bayesiano ingenuo con probabilità a priori
$0.15$:
\begin{equation}
\label{eq:osint}
O = \frac{0.15\,\ell_T}{0.15\,\ell_T + 0.85\,\ell_C},
\end{equation}
\[
\ell_T = \beta_T(b)\,\max(g, 0.01)\,\max(1 - e^{-0.3 s}, 0.01), \quad
\ell_C = \beta_C(b)\,\max(1 - g, 0.01)\,\max(e^{-0.6 s}, 0.01),
\]
con $\beta_T = 0.75$ o $0.25$ e $\beta_C = 0.01$ o $0.99$ secondo che la targa
sia o non sia in lista. Le verosimiglianze sono funzioni scelte per la
modellazione, non derivate dalle distribuzioni di Tabella~\ref{tab:init_sim}.

\emph{Canale comportamentale.}
$\beta = \operatorname{clip}\big((\operatorname{Var}(x) + \operatorname{Var}(y))/20,\, 0,\, 1\big)$
sulle ultime $10$ posizioni, con l'aggiunta di $\mathcal{U}(0.10,\, 0.25)$ per
i bersagli; nullo con meno di tre posizioni.

\emph{Fusione.} $\tau = 0.45\, v + 0.35\, O + 0.20\, \beta$. La banda di
variabilità è $\mu \pm 1.96\,\sigma$, calcolata sugli ultimi $10$ punteggi e
limitata a $[0,1]$.

\emph{Decisione.} L'azione nominale segue le soglie di
Tabella~\ref{tab:azioni}. L'azione \textsc{engage} è declassata ad
\textsc{alert} se l'ampiezza della banda supera $0.40$, oppure se più di tre
civili si trovano a distanza di Manhattan non superiore a $5$ e $\tau < 0.85$.

\begin{table}[htbp]
\centering
\begin{tabular}{llp{6.4cm}}
\toprule
\textbf{Azione} & \textbf{Condizione} & \textbf{Significato} \\
\midrule
\textsc{ignore} & $\tau < 0.22$ & entità non tracciata \\
\textsc{track} & $0.22 \le \tau < 0.38$ & entità sotto osservazione \\
\textsc{alert} & $0.38 \le \tau < 0.58$ & segnalazione a supervisione umana \\
\textsc{engage} & $\tau \ge 0.58$ & ingaggio \\
\bottomrule
\end{tabular}
\caption{Azioni e soglie sul punteggio di minaccia.}
\label{tab:azioni}
\end{table}

% figura automa: dopo conferma impostazione

\textbf{Risultati preliminari.} Con i difetti d'implementazione corretti, su
$150$ passi e venti repliche con seme diverso: isolando il solo canale visivo,
la frazione di osservazioni di bersaglio classificate \textsc{alert} o
\textsc{engage} è $0.506 \pm 0.017$ senza contromisura e $0.219 \pm 0.015$ con
contromisura, contro $R_1^{\text{pre}} = 0.501$ e $R_1^{\text{post}} = 0.220$
misurati; il tasso di falso allarme sui civili è $0.029$, contro
$1 - R_2 = 0.030$. Nel sistema a canali fusi la contromisura isolata non altera
la classificazione dei bersagli ($0.998$ contro $0.997$): il canale OSINT è
sufficiente a sostenerla. Combinata con l'attacco al canale OSINT, la
contromisura la riduce da $0.340$ a $0.150$, rapporto $0.44$ pari a
$R_1^{\text{post}}/R_1^{\text{pre}}$. Pesi, soglie e parametri delle
distribuzioni sono scelte di modellazione, non stime: il livello ha valore
dimostrativo del meccanismo di propagazione.


% C — Sez. 4.2: ipotesi white-box
\section{Modello di minaccia}
\label{sec:minaccia}

\textbf{Ipotesi di conoscenza del rilevatore.}
Si assume che l'avversario conosca l'architettura e i pesi del rilevatore
impiegato. L'assunzione è deliberatamente favorevole all'attaccante e
circoscrive la portata dei risultati: le misure riportate nel
Capitolo~5 costituiscono un limite superiore all'efficacia
conseguibile, non una stima operativa contro un rilevatore ignoto. La
trasferibilità della perturbazione verso un modello diverso non è stata
misurata ed è indicata fra gli sviluppi futuri.


% D — note puntuali: R2, NMS, split, valore atteso
\section{Chiarimenti su metriche e dati}
\label{sec:chiarimenti}

\textbf{Metriche.} $R_1$ misura l'efficacia della perturbazione sul bersaglio;
$R_2$ controlla che la perturbazione non introduca falsi allarmi nei
fotogrammi privi di bersaglio. La loro media geometrica $\sqrt{R_1 R_2}$ è la
metrica primaria. La soppressione dei non-massimi (\textit{Non-Maximum
Suppression}, NMS) elimina i riquadri sovrapposti riferiti al medesimo
oggetto: la confidenza $\bar{c}$ dell'Equazione~\eqref{eq:conf_aggregata} è
calcolata prima della NMS, le metriche di valutazione dopo, alla soglia di
confidenza $0.50$.

\textbf{Dati.} Si adottano le suddivisioni ufficiali di addestramento e di
validazione. VisDrone: $531$ immagini di validazione, le $548$ ufficiali meno
le $17$ prive di annotazioni di persona. Okutama-Action: $54\,664$ fotogrammi
di addestramento e $14\,210$ di validazione, sottocampionati con passo $27$
($n = 527$) per ridurre la correlazione fra fotogrammi consecutivi.
% TODO: n. immagini addestramento VisDrone

\textbf{Programma del tasso di apprendimento.} Il ``valore atteso'' citato nel
testo va inteso come valore teorico previsto dal programma a coseno, non come
aspettazione statistica.


% E — domande per la call
\section*{Domande aperte}

\begin{enumerate}
\item \textbf{Terminologia del livello di simulazione.} Lo spazio degli stati
non è finito. Va bene ``macchina a stati'' con la precisazione della
Sezione~\ref{sec:automa}, oppure preferisce un'altra dizione?

\item \textbf{Figura dell'automa.} (a) schema a blocchi stato $\to$
transizione $\to$ uscita; (b) diagramma della sola catena decisionale di un
passo; (c) entrambe.

\item \textbf{Collocazione delle tabelle di parametri} (EoT, ottimizzazione,
inizializzazione): (a) nel testo; (b) in appendice.

\item \textbf{Formulazione a soglia della funzione obiettivo:} la tolgo, come
suggerito nel suo commento?

\item \textbf{Risultati del livello di simulazione:} (a) tabella quantitativa
nel Capitolo 5; (b) solo forma descrittiva, come nel paragrafo sopra.

\item \textbf{Fusione.} Con probabilità a priori $0.50$ l'aggiornamento
bayesiano del punteggio fuso restituisce il valore in ingresso: la fusione è
quindi una combinazione lineare pesata, e la descrivo come tale. Concorda?
\end{enumerate}

\end{document}
